\section{Dot Product and Duality}
\label{sec:dot-product}

\noindent\textbf{Subject:} Linear Algebra for ML (3Blue1Brown) \\
\textbf{Status:} growing

\subsection*{Summary}
The dot product measures how aligned two vectors are. Duality connects the dot product to linear transformations from space to the number line.

\subsection{Dot Product}

\begin{defbox}
When you multiply two vectors with the same dimension, each row is multiplied by the corresponding row of the other, then the results are added together.
\end{defbox}

\[
\begin{bmatrix} 1 \\ 2 \end{bmatrix} \cdot \begin{bmatrix} 3 \\ 4 \end{bmatrix}
= 1 \times 3 + 2 \times 4 = 11
\quad \text{(Projection of a vector)}
\]

\paragraph{Interpretation:}
\begin{itemize}[nosep]
  \item If the dot product is \textbf{positive} $\rightarrow$ vectors are in the \textbf{same direction}
  \item If the dot product is \textbf{0} $\rightarrow$ vectors are \textbf{perpendicular}
  \item If the dot product is \textbf{negative} $\rightarrow$ vectors are in \textbf{opposite directions}
\end{itemize}

\[
\vec{v} = \begin{bmatrix} 4 \\ 1 \end{bmatrix}
\]

A \textbf{$1\times2$ matrix} $\leftrightarrow$ \textbf{2D vectors} $\rightarrow$ Matrix-vector product $\Leftrightarrow$ Dot product

\subsection{Duality}

\begin{defbox}
Duality of a vector is the \textbf{linear transformation that it encodes} --- a unique vector in which the linear transformation goes from space to the number line.
\end{defbox}

\subsection*{Related Topics}
\begin{itemize}[nosep]
  \item Section~\ref{sec:non-square} --- Non-Square Matrix
  \item Section~\ref{sec:cross-product} --- Cross Product
\end{itemize}

\bigskip
\noindent\textit{Tags: linear-algebra, dot-product, duality, math, phase-1}
